\section{Interpretation}

\paragraph{Why repayment status dominates.}
Next month's default is largely the next step of a delinquency process, so the latest status is nearly a lagged copy of the label. Arrears persist or deepen: 59.1\% of clients two months late in September, and 85.9\% of those three months late, were two months late in August (training rows; Table~\ref{tab:a-paycodes-rates}). Repayment is also the only liquidity signal the bank sees (income and debts are unseen). Adding September's status, bill and payment to limit and demographics lifts XGBoost's CV AUC from 0.63 to 0.77. Older months add 0.013 as raw months, or 0.017 through six-month summaries such as months in arrears, which tell a new slip from a chronic one (5/5 folds each; training CV, post-freeze; Table~\ref{tab:a-ablation-history}, Figure~\ref{fig:a-ablation}). A crisis-vintage model thus learns crisis roll rates (how fast arrears deepen) and may overstate them in calmer years (\S7, iii).

\paragraph{Where the ensembles' edge comes from: shape, then small interactions.}
Most of XGBoost's edge over the logit is one-variable curves; interactions add little. A depth-1 booster (one split per tree) adds one-variable effects of any shape; depth 3 lets variables interact (trees re-tuned per depth). The additive booster recovers 0.0225 of the 0.0261 gap (5/5 folds) and depth 3 adds only 0.0036, an upper estimate (selected on the same folds; training CV, post-freeze; Table~\ref{tab:a-ablation-depth}; same split on raw columns). The one-hot logit gives each status code its own risk and still trails, so the remaining shape likely sits in other columns, such as utilisation's U in partial-dependence plots (PDPs [10]); we did not measure which. Friedman's $H$ (the share of a pair's effect beyond the two separate effects) is 0.061 for status $\times$ limit: little support for strong interactions in the pairs tested (Figures~\ref{fig:pdp-main}, \ref{fig:pdp-interaction}). A binned scorecard or additive model should thus recover most of the gain and stay explainable.

\paragraph{Which variables matter, and in which direction.}
All four measures put September status first, then disagree for known reasons. The forest's impurity importance (MDI) credits every training-row split. Each split sees only a random subset of columns, so when September status is not on offer, correlated delay columns stand in and collect its credit. Permutation, shuffling a column on rows a tree did not see (out of bag) [3], shows how much \emph{this fitted model} leans on it. There the August and July codes, 4th and 5th on MDI, rank only 22nd and 24th. Refitting without a column measures what is lost: of six columns refitted, only September status costs over 0.001 out-of-bag AUC, and the delay counters recover most of even that (Table~\ref{tab:a-refit}). Across models, SHAP (XGBoost) and permutation (forest) agree (rank correlation 0.77), MDI and permutation far less (Table~\ref{tab:a-importance}); September status, worst delay and months in arrears carry one signal, so read them as a block. A Shapley value [14] splits a client's score into feature contributions summing to its log-odds minus the average: arrears push towards default and a larger limit away (Figure~\ref{fig:shap-beeswarm}), an association (\S2). These are reasons a bank can give a customer [15]; they describe how the fitted model behaves, not the borrowers.

\paragraph{Monotone constraints: impose only what credit logic guarantees.}
A monotone constraint fixes a direction in advance: all else equal, the score never falls as a client falls further behind. The frozen version also forced it up from $-2$ to $-1$ to 0 and with utilisation and balance trend (Table~\ref{tab:a-monotone}). \S2 contradicts both the order and the rise, so this costs 0.0019 CV AUC (no fold better). Credit logic is certain only from revolving credit upwards. Up to three months late the free booster already obeys: moving 200 training clients from revolving credit to one, two and three months late (counters updated) never lowers a score. Its wrong-way steps lie among non-delinquent codes or in the thin 4- and 5-month cells (57 and 16 clients), where observed default also falls. Constraining only delay (score up) and limit (score down) is within 0.0005 of the free booster (training CV, post-freeze; Table~\ref{tab:a-monotone-respec}). A bank can thus promise, almost for free, that slipping from revolving credit into arrears never lowers a score (rare jumps from $-2$ or $-1$ into arrears still lower it slightly).

\paragraph{Personal attributes: little ranking power, real fairness trade-offs.}
Sex, age, marriage and education barely help ranking: removing them costs XGBoost 0.002 CV AUC (no fold better without; training CV, post-freeze; Table~\ref{tab:a-ablation-demo}). Yet its cross-fitted 5:1 cost rises in all five folds, as age carries information the limit does not. When groups default at different rates, no score can be both calibrated within each group and equal in its error rates [7, 13]. Men default 1.19 times as often as women. XGBoost's PDs match each large group's observed rate, so it flags men 1.18 times as often, partly through SEX and MARRIAGE; men who would repay are therefore flagged more often too (training OOF, post-freeze; Table~\ref{tab:a-fairness}). The logit flags men more often still, as one SEX coefficient shifts every man's log-odds alike and more of its scores sit near the cut-off. Dropping the columns narrows the sex gap but under-predicts men, flags more under-30s (without age, their low limits make them look riskier) and leaves proxies, mainly the limit (Figure~\ref{fig:a-fairness}). ``Fairness through unawareness'' is thus a trade-off a bank must choose and document, not a fix; present-day frameworks class such scoring as high-risk [8] and want attributes justified and reviewed [15].
